% SHARD-ID: CA-60-QUBIT-CANDIDATE-SCAN-RESULTS
% SHARD-TITLE: First Candidate Scan Results
% SHARD-SUMMARY: Reports the checked 99-point qubit Hamiltonian scan.
% SHARD-SUMMARY: All points fail before a nontrivial exact boost witness; TFIM and XXZ fail at that boost gate.
% SHARD-SUMMARY: The result rejects the exact first-moment route, not the sourced CFT status of known critical models.
% SHARD-KEYWORDS: scan results, TFIM, XXZ, Heisenberg, boost witness, route failure

\section{First Candidate Scan Results}
\label{sec:qubit-candidate-scan-results}

\subsection{Status, Sources, and Dependencies}

\textbf{Status: checked run result.}  This shard reports only the run bundle
named in CA-59.

Dependencies: CA-53--CA-59.

% Run: runs/2026-05-31-qubit-candidate-scan/summary.toml.
% Source boundary: CA-55.

\subsection{Headline Table}

The 99 samples split as:
\[
\begin{array}{c|c}
\text{verdict} & \text{count}\\
\hline
\texttt{excluded\_current\_collapsed} & 9\\
\texttt{excluded\_no\_conservation\_witness} & 27\\
\texttt{excluded\_no\_boost\_witness} & 63
\end{array}
\]
No point reaches \texttt{not\_excluded\_algebraic}.

\subsection{Notable Points}

The self-dual TFIM sample
\[
  -ZZ-\frac12(XI+IX)
\]
is currentful and has a conservation witness, but its exact boost residual norm
is nonzero in the run.  The Heisenberg point and the \(XX\) point behave the
same way: they reach the boost gate and fail it.

Because TFIM and XXZ have local source anchors as critical or gapless models,
the conclusion is about the route: the exact CA-39 boost witness is too strict
as a universal candidate filter.
